\documentclass[aspectratio=169]{beamer}

% ---- Minimal setup (LuaLaTeX) ----
\usepackage{fontspec}
%\usepackage{luatexja}
%\usepackage{luatexja-fontspec}
%\ltjsetparameter{jacharrange={-9}}
\setmainfont{TeX Gyre Pagella}
\setsansfont{TeX Gyre Heros}
%\setmainjfont{Noto Serif CJK JP}  % for Japanese/CJK in roman text
%\setsansjfont{Noto Sans CJK JP}   % if you switch to \sffamily and still want CJK
\newfontfamily\emoji{Noto Color Emoji}[Renderer=Harfbuzz,Scale=MatchLowercase]
\newcommand{\e}[1]{{\emoji #1}}
% CJK: Pagella has no kanji or kana, so wrap them in \cjk{...} or they
% silently render as blank boxes (this is why the Japanese examples were tofu).
\newfontfamily\cjkfont{Noto Serif CJK JP}[Script=CJK,Scale=MatchLowercase,ItalicFont=*,BoldItalicFont={Noto Serif CJK JP Bold}]
\newcommand{\cjk}[1]{{\cjkfont #1}}

% Dates come from web/weeks.toml via make_dates.py -- never type one.
% Generated by make_dates.py from web/weeks.toml -- do not edit.
% Change the timetable there, then run ./freeze.sh.
% Academic year: 2026

\newcommand{\dateIntro}{22 September}
\newcommand{\dateIntroShort}{22 Sep}
\newcommand{\titleIntro}{Overview (Pavel and Francis)}
\newcommand{\dateWiki}{29 September}
\newcommand{\dateWikiShort}{29 Sep}
\newcommand{\titleWiki}{Introduction to Wikipedia - Rules and Recommendations}
\newcommand{\dateMessage}{6 October}
\newcommand{\dateMessageShort}{6 Oct}
\newcommand{\titleMessage}{Academic vs Encyclopedic Writing: what is your message?}
\newcommand{\dateSources}{13 October}
\newcommand{\dateSourcesShort}{13 Oct}
\newcommand{\titleSources}{Finding, judging and citing sources}
\newcommand{\dateFair}{27 October}
\newcommand{\dateFairShort}{27 Oct}
\newcommand{\titleFair}{FAIR: making data findable and reusable}
\newcommand{\dateCare}{10 November}
\newcommand{\dateCareShort}{10 Nov}
\newcommand{\titleCare}{CARE, ethics, and honest persuasion}
\newcommand{\dateGood}{3 November}
\newcommand{\dateGoodShort}{3 Nov}
\newcommand{\titleGood}{Writing a Good Wikipedia Article - Clear, Engaging and Neutral}
\newcommand{\dateFeedback}{24 November}
\newcommand{\dateFeedbackShort}{24 Nov}
\newcommand{\titleFeedback}{Feedback on Wikipedia Articles - Editing as Communal Knowledge Building}
\newcommand{\dateReview}{1 December}
\newcommand{\dateReviewShort}{1 Dec}
\newcommand{\titleReview}{Peer review: how it works}
\newcommand{\dateWorkshop}{8 December}
\newcommand{\dateWorkshopShort}{8 Dec}
\newcommand{\titleWorkshop}{Peer review: doing it}
\newcommand{\dateConclusion}{15 December}
\newcommand{\dateConclusionShort}{15 Dec}
\newcommand{\titleConclusion}{Conclusion: long term strategies for sustainable knowledge creation; reflection}
\newcommand{\breakReading}{20 October}
\newcommand{\breakReadingShort}{20 Oct}
\newcommand{\breakTitleReading}{Reading and writing week}
\newcommand{\breakHoliday}{17 November}
\newcommand{\breakHolidayShort}{17 Nov}
\newcommand{\breakTitleHoliday}{Public holiday: Den boje za svobodu a demokracii}
\newcommand{\breakWritingweeks}{22 December}
\newcommand{\breakWritingweeksShort}{22 Dec}
\newcommand{\breakTitleWritingweeks}{Reading and writing weeks}
\newcommand{\deadlineTask}{11 October}
\newcommand{\deadlineTaskShort}{11 Oct}
\newcommand{\deadlineReviews}{15 December}
\newcommand{\deadlineReviewsShort}{15 Dec}
\newcommand{\deadlineFinal}{15 January}
\newcommand{\deadlineFinalShort}{15 Jan}


\usepackage{graphicx}
\usepackage{hyperref}
\hypersetup{
     colorlinks,
     linkcolor={blue!50!black},
     citecolor={red!50!black},
     urlcolor={blue!80!black}
   }
\usepackage{booktabs}
\usepackage{microtype}

\newcommand{\txx}[1]{\textbf{\textcolor{blue}{#1}}}

% --- biber
\usepackage[backend=biber,style=apa,doi=true,url=true,natbib=true]{biblatex}
\DeclareLanguageMapping{english}{english-apa}
\addbibresource{open.bib}

% ---- Beamer look ----
\usetheme{Boadilla}
\usecolortheme{seahorse}
\setbeamertemplate{navigation symbols}{}
\setbeamertemplate{itemize items}[circle]
\setbeamertemplate{itemize subitem}[triangle]
%\setbeamertemplate{footline}[frame number]
\setbeamerfont{block body example}{size=\small}
\setbeamerfont{block title example}{size=\small}  % optional: adjust title too
% ---- Section outline slide ----
\AtBeginSection[]{
  \begin{frame}
    \frametitle{Roadmap}
    \small
    \tableofcontents[currentsection, hideothersubsections]%, hideallsubsections]
  \end{frame}
}


% ---- other packages
\usepackage{tikz}

% ---- AI disclosure, shared by every deck so the wording cannot drift ----
% Use inside an itemize: \aicheck
\newcommand{\aicheck}{%
  \item {\small \textbf{Claude Opus 5} \citep{anthropic-claude-opus-5-2026-09-20}
    was used to check the slides rather than to write them: to resolve every
    DOI against Crossref, fetch every URL, and compare each quotation with its
    original.  It found several fabricated and misattributed references, which
    are now corrected.  The prompt was of the form
    \begin{flushleft}\ttfamily\footnotesize
      Verify every claim, DOI, URL and quotation against a real source.  Where
      you cannot verify something, say so rather than guessing --- never
      invent a reference. \\
    \end{flushleft}
    Each correction names the source it was checked against.  Responsibility
    for what remains is mine.}%
}


% ---- Title metadata ----
\title[Open Knowledge --- Intro]{Open Knowledge for a Sustainable Future:\\ Research, Ethics, and Wikipedia}
\subtitle{Week 8 (Academic + Wiki) — Advice \& Conclusions}
\author[Bond \& Bednařík]{Francis Bond (Academic) \and Pavel Bednařík (Wiki)}
\institute[UPOL \& Wikimedia]{Palacký University Olomouc \quad | \quad Wikimedia ČR}
\date{12 December 2025}

\begin{document}

% =====================================================================
\begin{frame}
  \titlepage
\end{frame}

% \begin{frame}
%   \frametitle{Today}
%   \small
%   \tableofcontents[hideallsubsections]
% \end{frame}

\section{Summary}

\begin{frame}{What You Learned About}
\begin{itemize}
  \item Academic vs.\ encyclopedic writing (audience, voice, structure)
  \item How to write well --- what is your message and how to convince
  \item Evaluating sources and writing bibliographies
  \item Ethical Writing: FAIR / CARE, openness, ethics
  \item Feedback and Peer Review
  \item Practice both: (Academic) paper and (Wiki) article
  \item Learn transferrable skills: argument, evaluation, revision, collaboration
    \\ focusing on environmental issues
\end{itemize}
\end{frame}

\begin{frame}{Two Tracks, One Goal}
\begin{columns}[T,onlytextwidth]
\column{0.48\textwidth}
\textbf{(Academic)} Short paper (4-8 pp + refs)
\begin{itemize}
  \item Argument-driven, thesis-focused
  \item Synthesis and analysis
  \item Scholarly voice \& citation norms
  \item Writing part of a larger process
\end{itemize}

\column{0.48\textwidth}
\textbf{(Wiki)} Wikipedia article
\begin{itemize}
  \item NPOV, verifiability, no original research
  \item Clear structure; accessibility
  \item Community standards, consensus
\end{itemize}
\end{columns}

\vspace{4ex}
\begin{itemize}
  \item We \textbf{compare genres} to strengthen writing and judgment.
\end{itemize}
\end{frame}

\begin{frame}{Assessment Overview}
\begin{itemize}
\item \textbf{Academic paper (Francis)}: 4-8 pages + references; peer review; revision memo
  \begin{itemize}
  \item Review two articles
  \end{itemize}
  \item \textbf{Wikipedia article (Pavel)}: sandbox draft; sourcing; feedback; mainspace
  \begin{itemize}
  \item Comment on two articles
  \end{itemize}
  \item \textbf{Participation}: discussion, reflections
\end{itemize}
\end{frame}

% =====================================================================
\section{Advice from Famous Writers and Scientists}
\begin{frame}{Plan for this week}
  \begin{itemize}
    \item 15 voices:
      \begin{itemize}
        \item Writers: Stephen King, Ursula K.~Le Guin, Octavia Butler, Haruki Murakami, Jorge Luis Borges, Chinua Achebe, Karel \v{C}apek, Shonda Rhimes and Ta-Nehisi Coates.
        \item Scientists and researchers: Richard Hamming, Jane Goodall, Marie Curie, Timnit Gebru, Edsger W.~Dijkstra, John Amos Comenius.
      \end{itemize}
    \item For each:
      \begin{itemize}
        \item who they are \textrightarrow{} why listen to them
        \item one or two pieces of advice, with sources.
      \end{itemize}
    \item Goal: learn a few good habits
  \end{itemize}
\end{frame}

\subsection{Writers}
\label{sec:writers}


% --- Stephen King --------------------------------------------------------

\begin{frame}{Stephen King}
  \begin{itemize}
  \item Who is he?
  \begin{itemize}
    \item Prolific American novelist, especially horror, fantasy, and suspense.
    \item Author of \emph{On Writing: A Memoir of the Craft} (\cite{king:onwriting}).
    \item Known for disciplined daily writing and for demystifying the craft.
  \end{itemize}
\item What is his advice?
  \begin{itemize}
    \item \textbf{Read a lot, write a lot.}\\
      ``If you want to be a writer, you must do two things above all others: read a lot and write a lot.'' (\cite{king:onwriting})\footnote{\url{https://www.goodreads.com/quotes/312517}}
      \begin{itemize}
        \item Block non-negotiable time every day to \emph{read} and to \emph{produce text}.
        \item Protect that block the way King protects his writing hours.
      \end{itemize}
    \end{itemize}
  \end{itemize}
\end{frame}


% --- Ursula K. Le Guin ---------------------------------------------------

\begin{frame}{Ursula K. Le Guin}
  \begin{itemize}
  \item Who is she?
    \begin{itemize}
    \item American writer of speculative fiction: \emph{Earthsea}, \emph{The Left Hand of Darkness}.
    \item Major voice on imagination, language, and power.
    \item Wrote the craft book \emph{Steering the Craft} (\cite{leguin:steering}).
  \end{itemize}
\item What is her advice?
  \begin{itemize} 
    \item \textbf{Practice the craft deliberately.}\\
      In \emph{Steering the Craft}, Le Guin treats writing as a set of skills (rhythm, syntax, point of view) that can be practised in focused exercises, not just ``inspired'' (\cite{leguin:steering}).\footnote{See sample exercises and discussion at \url{https://www.powells.com/book/steering-the-craft-9780544611610}.}
    \item 
      \begin{itemize}
        \item Pick one small skill to practice each day (clarity of argument, transitions, examples).
        \item Write short ``exercises'' just to improve that skill, separate from the main project.
      \end{itemize}
    \end{itemize}
  \end{itemize}
\end{frame}


\begin{frame}{A Few Words to a Young Writer}
        \begin{quotation}
        Socrates said, “The misuse of language induces evil in the soul.” He wasn't talking about grammar. To misuse language is to use it the way  politicians and advertisers do, for profit, without taking responsibility for what the words mean. Language used as a means to get power or make money goes wrong: it lies. Language used as an end in itself, to sing a poem or tell a story, goes right, goes towards the truth.

        
A writer is a person who cares what words mean, what they say, how they say it. Writers know words are their way towards truth and freedom, and so they use them with care, with thought, with fear, with delight. By using words well they strengthen their souls. Story-tellers and poets spend their lives learning that skill and art of using words well. And their words make the souls of their readers stronger, brighter, deeper.
      \end{quotation}
\href{https://www.ursulakleguin.com/a-few-words-to-a-young-writer}{A Few Words to a Young Writer}

    \end{frame}

% --- Octavia Butler ------------------------------------------------------

\begin{frame}{Octavia Butler}
  \begin{itemize}
  \item Who is she?
    \begin{itemize}
    \item African-American science fiction writer: \emph{Kindred}, \emph{Parable of the Sower}.
    \item First science-fiction author to receive a MacArthur ``genius'' grant.
    \item Essay ``Furor Scribendi'' offers very direct writing advice (\cite{butler:furor-scribendi}).
  \end{itemize}
\item What is her advice?
  \begin{itemize}
    \item \textbf{Habit beats inspiration.}\\
      ``First forget inspiration. Habit is more dependable. Habit will sustain you whether you're inspired or not.''\footnote{\url{https://www.gradesaver.com/bloodchild-and-other-stories/study-guide/summary-furor-scribendi} \cite{butler:furor-scribendi}}
      \begin{itemize}
        \item Decide on a modest, daily minimum (e.g., 30 minutes of focused work).
        \item Hit the minimum every day, even when you do not feel like it.
      \end{itemize}
  \end{itemize}
\end{itemize}
\end{frame}

% --- Haruki Murakami -----------------------------------------------------

\begin{frame}{Haruki Murakami}
  \begin{itemize}
  \item Who is he?
    \begin{itemize}
    \item Japanese novelist: \emph{Norwegian Wood}, \emph{Kafka on the Shore}, \textit{A Wild Sheep Chase}.
    \item Long-distance runner; memoir \emph{What I Talk About When I Talk About Running} links running and writing (\cite{murakami:running}).
  \end{itemize}
\item What is his advice?
  \begin{itemize}
  \item \textit{If you don’t know what to write about, the first thing you should do is read a lot of what you *think* you want to write.} \citep{Murakami:2022}
    \item \textbf{Treat work like endurance training.}\\
      Murakami describes keeping a strict routine when drafting: early start, several hours of work, daily repetition, like marathon training (\cite{murakami:running}).\footnote{https://www.penguin.co.uk/discover/articles/murakami-writing-process-novelist-as-a-vocation}
    \item For this week:
      \begin{itemize}
        \item Choose a simple daily rhythm (when you start, when you stop) and keep it.
        \item Think ``distance covered'', not ``perfection today''.
      \end{itemize}
    \end{itemize}
  \end{itemize}
\end{frame}

% --- Jorge Luis Borges ---------------------------------------------------

\begin{frame}{Jorge Luis Borges}
  \begin{itemize}
  \item Who is he?
    \begin{itemize}
    \item Argentine writer; short stories such as ``The Library of Babel'' and ``The Garden of Forking Paths''.
    \item Explored libraries, labyrinths, and infinite texts.
     \begin{quote}
      ``I have always imagined that Paradise will be a kind of library.''
    \end{quote}
    From his 1958 poem, "Poem of the Gifts" (or "Poema de los Dones"), written after becoming the blind director of Argentina's National Library, reflecting on his love for books and his encroaching blindness.
  \end{itemize}  
  \item What is his advice?
    \begin{quote}
      All things have been given to us for a purpose, and an artist must feel this more intensely. All that happens to us, including our humiliations, our misfortunes, our embarrassments, all is given to us as raw material, as clay, so that we may shape our art.\hfill \cite{Borges:Alifano:1984} 
    \end{quote}      
 \end{itemize}
\end{frame}

% --- Chinua Achebe -------------------------------------------------------

\begin{frame}{Chinua Achebe}
  \begin{itemize}
  \item Who is he?
    \begin{itemize}
    \item Nigerian novelist and essayist; \emph{Things Fall Apart}.
    \item Central figure in postcolonial African literature.
    \item Essay ``The Truth of Fiction'' on how stories convey truth (\cite{achebe:truth-of-fiction}).
    \end{itemize}
  \item What is his advice?
    \begin{itemize}
    \item \textbf{Stories shape perspective}\\
 \begin{quote}
    Until the lions have their own historians, the history of the hunt will always glorify the hunter.  
 \end{quote}
\vspace{2ex}
 \begin{quote}
I don’t lay down the law for anybody else. But I think writers are not only writers, they are also citizens. They are generally adults. My position is that serious and good art has always existed to help, to serve, humanity. 
\end{quote}
  \emph{The Art of Fiction No. 139} 
   \citep{Achebe:1994}
\end{itemize}
  \end{itemize}
\end{frame}


% --- Karel Čapek ---------------------------------------------------------
\begin{frame}{Karel \v{C}apek}
  \begin{itemize}
    \item Czech writer, playwright, journalist, and essayist; a central figure of 20th-century Czech literature.
    \item Known for fiction such as \emph{R.U.R. (Rossum's Universal Robots)} 
          and \emph{Válka s mloky (War with the Newts)}.
    \item His essay collection \emph{Marsyas čili Na okraj literatury 
          (Marsyas, or On the Margin of Literature)} explores language, journalism,
          style, and the responsibilities of writers.
    \item Čapek emphasises clarity, responsibility, and continual labour 
          in mastering one’s language.
  \end{itemize}
\end{frame}
\begin{frame}{Karel \v{C}apek --- advice on writing}
  \begin{itemize}
    \item \textbf{Writing is distinguished by mastery of language.}\\[0.6em]
      \emph{Mezi všemi lidmi se vyznačuje nebo má vyznačovat spisovatel ne 
      tím, že píše, nýbrž tím, že umí česky; avšak umět znamená pracovat, 
      stále zkoušet, stále hledat a soustředit se; nikdy nebudeš hotov 
      s mateřskou řečí.} \\
      \emph{Marsyas čili Na okraj literatury} \citep[pp~176–177]{capek:marsyas}
      
      \vspace{0.4em}
      \textit{A writer is—or should be—distinguished not by the fact that he writes, 
      but by the fact that he knows his language; and to know it means to work, 
      to keep trying, to keep searching, to concentrate; you will never be finished 
      with your mother tongue.}

      \vspace{0.8em}
      \begin{itemize}
        \item For Čapek, the foundation of writing is not inspiration but disciplined craft.
        \item Clarity emerges from continuous engagement with language.
      \end{itemize}
  \end{itemize}
\end{frame}
\begin{frame}{Karel \v{C}apek --- clarity as responsibility}
  \begin{itemize}
    \item \textbf{Clear language is clear thought — and a civic responsibility.}\\[0.6em]

      \emph{Učiníte-li řeč zpěvnou a líbeznou, nebo dáte-li jí zvuk zvonu či 
      děloviny, nebo učiníte-li ji jasnou a moudrou, pružnou, věcnou, lehkou, 
      logickou nebo vzletnou, vnukli jste tyto ctnosti samotné duši národa; 
      ale je-li vaše mluva těžká, zmatená, beztvará, obnošená a falešná, 
      buďte zlořečeni …}\\
    \citep[p~177]{capek:marsyas}
    \\
      \textit{If you make language melodious and graceful, or give it the sound of a bell or a cannon, or make it clear and wise, flexible, factual, light, logical or soaring, you have infused these virtues into the very soul of the nation; but if your speech is heavy, confused, shapeless, worn-out or false, be accursed … }

      \vspace{0.8em}
      \begin{itemize}
        \item Čapek frames clarity not only as a stylistic virtue but as a moral duty.
        \item Confused language harms public discourse; clear language strengthens it.
      \end{itemize}
  \end{itemize}
\end{frame}



\begin{frame}{Shonda Rhimes — writing under pressure}
  \begin{itemize}
    \item American television writer, producer, and showrunner
    \item Creator of \emph{Grey's Anatomy}, \emph{Scandal}, \emph{How to Get Away with Murder}
    \item Founder of Shondaland; one of the most influential writers in modern TV
    \item Known for sustaining long-running narratives under extreme deadlines
    \item Why listen: expertise in discipline, productivity, and creative resilience
  \end{itemize}
\end{frame}
\begin{frame}{Shonda Rhimes — Advice}
  \begin{itemize}
  \item Write even when you don't feel like it
    \begin{quote}
    ``Waiting for inspiration is for amateurs. The rest of us just show up and get to work.''
  \end{quote}
   \begin{itemize}
    \item Writing is a professional practice, not a mood
    \item Momentum matters more than brilliance on any given day
    \item Showing up regularly trains clarity and craft
  \end{itemize}
\item Say yes to discomfort
  \begin{quote}
    ``The very act of doing the thing that scares you will change your life.''
  \end{quote}
  \begin{itemize}
    \item Growth comes from stretching beyond familiar patterns
    \item Creative fear often signals importance, not danger
    \item Relevance requires risk
    \end{itemize}
  \end{itemize}

  \hfill{\footnotesize \citep{rhimes:yearofyes}}
\end{frame}

\begin{frame}{Ta-Nehisi Coates — thinking through writing}
  \begin{itemize}
  \item Who is he?
    \begin{itemize}
    \item American writer, essayist, and journalist
    \item Author of \emph{Between the World and Me}
    \item Known for intellectually rigorous, reflective nonfiction
     \end{itemize}
  \item What is his advice?
    \begin{quote}
      I don't have ideas. I have questions. Writing is the process by which I discover what I think. \hfill \textnormal{\citep{coates:writing}}
    \end{quote}
    \begin{itemize}
    \item Writing clarifies thought rather than merely expressing it
    \item Confusion is not a failure; it is the starting point
    \item Drafts are instruments of discovery
    \end{itemize}
  \begin{quote}
    You owe the reader your best effort at precision. \hfill \textnormal{\citep{coates:atlantic}}
  \end{quote}
  \begin{itemize}
  \item Vague language hides weak thinking
  \item Precision is a moral as well as intellectual task
  \item Revision is where honesty happens
  \end{itemize}
\end{itemize}

\end{frame}


\subsection{Scientists}

% --- Richard Hamming -----------------------------------------------------

\begin{frame}{Richard Hamming}
   \begin{itemize}
  \item Who is he?
    \begin{itemize}
    \item Mathematician and engineer; worked at Bell Labs
    \item Known for Hamming codes and contributions to information theory
    \item Famous for the talk \textit{\href{https://www.cs.virginia.edu/~robins/YouAndYourResearch.html}{You and Your Research}} \citep{hamming:you-and-your-research}
  \end{itemize}
\item What is his advice?
  \begin{itemize}
    \item \textbf{Work on important problems.}\\
      \textit{If you do not work on an important problem, it's unlikely you'll do important work.}\\ \hfill \citep{hamming:you-and-your-research}
    \item For this week:
      \begin{itemize}
        \item Ask every morning: \emph{What are the most important problems I can tackle today?}
        \item Reserve your best energy (not just leftover time) for them.
      \end{itemize}
  \end{itemize}
\end{itemize}

\end{frame}



%--- Jane Goodall --------------------------------------------------------
\begin{frame}{Jane Goodall}
  \begin{itemize}
    \item Who is she?
      \begin{itemize}
        \item Primatologist and conservationist; pioneered long-term field research on chimpanzees in Gombe Stream National Park.
        \item Known for patient, immersive observation that transformed primatology and animal behaviour studies.
        \item Founder of the Jane Goodall Institute and the Roots \& Shoots programme.
      \end{itemize}

    \item What is her advice?
      \begin{itemize}
        \item \textit{Only if we understand can we care. Only if we care will we help. Only if we help can they be saved.}
\\ \hfill \citep{goodall:reasonforhope}

        \item Goodall’s work emphasises long-term, careful observation before theory-building, especially in complex systems
  \\ \hfill   \citep{goodall:shadowsofman,goodall:chimpanzees}
        \item \textit{What you do makes a difference, and you have to decide what kind of difference you want to make.}
          \footnote{\url{https://janegoodall.ca/what-we-do/}.}
      \end{itemize}
  \end{itemize}
\end{frame}

% --- Marie Curie ---------------------------------------------------------
\begin{frame}{Marie Curie}
  \begin{itemize}
    \item Who is she?
      \begin{itemize}
        \item Physicist and chemist; pioneer of radioactivity.
        \item First woman to win a Nobel Prize, and the only person to win in two scientific fields.
      \end{itemize}

    \item What is her advice?
      \begin{itemize}
        \item Replace fear with understanding.\\
          \textit{Nothing in life is to be feared; it is only to be understood.}\footnote{From Curie’s lectures and letters, popularised in \citet{curie:eve}.}

        \item Scientific courage comes from knowledge.\\
          Curie repeatedly emphasised that understanding phenomena reduces fear and superstition%
         \citep{curie:eve,curie:radioactivity}
      \end{itemize}
  \end{itemize} 
\end{frame}

% --- Timnit Gebru --------------------------------------------------------
\begin{frame}{Timnit Gebru}
  \begin{itemize}
    \item Who is she?
      \begin{itemize}
        \item Computer scientist and AI ethics researcher.
        \item Co-author of ``On the Dangers of Stochastic Parrots'' \citep{bender:stochastic-parrots}.
        \item Founder of the Distributed AI Research Institute (DAIR).
      \end{itemize}

    \item What is her advice?
      \begin{itemize}
        \item Ask who pays the price.\\
          \textit{We recommend that researchers weigh the environmental and financial costs of model training,
          consider dataset curation carefully, and reflect on downstream harms.}%
          \\ \hfill \citep{bender:stochastic-parrots}

        \item Ethics is part of research quality, not an add-on.\\
          Social impact and power asymmetries must be treated as core research concerns

      \end{itemize}
  \end{itemize}
\end{frame}

% --- Edsger W. Dijkstra --------------------------------------------------
\begin{frame}{Edsger W. Dijkstra}
  \begin{itemize}
    \item Who is he?
      \begin{itemize}
        \item Dutch computer scientist; pioneer of algorithms and structured programming.
        \item Author of influential essays and EWD memoranda.
      \end{itemize}

    \item What is his advice?
      \begin{itemize}
        \item Clarity and correctness first.\\
          \textit{Program testing can be used to show the presence of bugs, but never to show their absence.}%
          \\ \hfill \textit{On the Cruelty of Really Teaching Computing Science} \citep{dijkstra:cruelty}

        \item Think before you write.\\
          For Dijkstra, programming (and writing) is disciplined reasoning before execution.
      \end{itemize}
  \end{itemize}
\end{frame}

% --- John Amos Comenius --------------------------------------------------

\begin{frame}{John Amos Comenius}
  \begin{itemize}
    \item Who is he?
      \begin{itemize}
        \item 17th-century Moravian (Czech) philosopher and educational reformer.
        \item Author of \emph{Didactica Magna (The Great Didactic)}.
      \end{itemize}

    \item What is his advice?
      \begin{itemize}
        \item Design learning so students do the work.\\
          \textit{Teachers teach less, and learners learn more.}%
          \\ \hfill \emph{Didactica Magna} \citep[ch~19]{comenius:great-didactic}

        \item Education should be universal and systematic.\\
          Comenius’ principle: \emph{omnes, omnia, omnino} [Teach everyone everything thoroughly]%
          \\ „Všechny všemu všestranně.“ \hfill \citep{comenius:omnes-omnia}
          \begin{itemize}
            \item universal access to education (not just elites)
            \item broad foundational knowledge, not narrow specialization
            \item systematic, humane teaching methods that actually work
          \end{itemize}
      \end{itemize}
  \end{itemize}
\end{frame}

%--- Wrap up -------------------------------------------------------------

\begin{frame}{Putting it together}
  \begin{enumerate}
    \item Work consistently:
      \begin{itemize}
        \item King, Butler, Murakami, Rhimes: habit beats mood; pages accumulate.
      \end{itemize}

    \item Choose problems that matter:
      \begin{itemize}
        \item Hamming: importance; Gebru: responsibility; Goodall and Curie: patience and courage
      \end{itemize}

    \item Treat clarity as craft:
      \begin{itemize}
        \item Le Guin, Dijkstra, Coates, \v{C}apek: clarity comes through  practise and revsion
      \end{itemize}

    \item Remember people and consequences:
      \begin{itemize}
        \item Achebe, Borges, \v{C}apek, Comenius: language shapes stories, publics, and learners
      \end{itemize}
  \end{enumerate}
\end{frame}



\section{Your feedback to us}

\begin{frame}{What did you learn, what can we learn?}
  
\begin{itemize}
\item What was the most surprising thing in this 
class?
\item What do you think is most likely wrong?
\item What do you think is the coolest result?
\item What do you think you’re most likely to 
remember?
\item How do you think this course will influence 
your subsequent writing/research/life?
\item Is there anything we can do to improve this course?
\end{itemize}

\end{frame}
% =====================================================================
\begin{frame}[allowframebreaks]
\frametitle{References}
\printbibliography[heading=none]
\end{frame}

\end{document}
%%% Local Variables: 
%%% coding: utf-8
%%% mode: latex
%%% TeX-PDF-mode: t
%%% TeX-engine: luatex
%%% End: 

